% !TEX root = ../main.tex
% Figure: Complete PCSEL design workflow from specs to fabrication
% Chapter: ch22 - Practical Simulation Workflow and Design Guide
% Description: Step-by-step flowchart covering entire development cycle

\begin{figure}[htbp]
  \centering
  \begin{tikzpicture}[scale=0.78, transform shape,
    step/.style={rectangle, draw, rounded corners, fill=#1!20, minimum width=3.5cm, minimum height=1cm, align=center},
    decision/.style={diamond, draw, aspect=1.8, fill=#1!20, minimum width=3cm, align=center},
    arrow/.style={->, >=stealth, thick},
    feedback/.style={->, >=stealth, thick, dashed}
  ]
    % Stage 1: Requirements
    \node[step=blue, minimum height=1.2cm] (specs) at (0,10) {Step 1: 明确技术指标\\波长$\lambda$, 功率$P$, 光束质量$M^2$, 工作温度};
    
    % Stage 2: Initial design
    \node[step=green, below=1cm of specs] (initial) {Step 2: 初始结构设计\\外延堆栈、PhC 晶格类型、预估$d_{etch}$};
    \draw[arrow] (specs) -- (initial);
    
    % Stage 3: Band structure calc
    \node[step=yellow, below=1cm of initial] (pwem) {Step 3: PWEM 能带计算\\寻找$\Gamma$点 flat band, 确定$a/\lambda$};
    \draw[arrow] (initial) -- (pwem);
    
    % Decision: Single mode?
    \node[decision=orange, below=1.5cm of pwem] (singlemode) {单模条件满足？\\$\Delta\nu > $增益线宽};
    \draw[arrow] (pwem) -- (singlemode);
    \draw[feedback] (singlemode.west) -- ++(-2,0) |- node[pos=0.7,left,align=center] {否\\调整$f$或$d_{etch}$} (initial.west);
    
    % Stage 4: Threshold analysis
    \node[step=cyan, below=1.5cm of singlemode] (threshold) {Step 4: 阈值增益分析\\CMT/FDTD 提取$\kappa$, 计算$g_{th}$};
    \draw[arrow] (singlemode) -- node[right,pos=0.5] {是} (threshold);
    
    % Decision: g_th achievable?
    \node[decision=purple, below=1.5cm of threshold] (gaincheck) {$g_{th} < g_{max}$?\\材料增益能力验证};
    \draw[arrow] (threshold) -- (gaincheck);
    \draw[feedback] (gaincheck.east) -- ++(2,0) |- node[pos=0.7,right,align=center] {否\\增强限制或减损耗} (initial.east);
    
    % Stage 5: Far-field optimization
    \node[step=pink, below=1.5cm of gaincheck] (farfield) {Step 5: 远场优化\\调节 Apodization, 抑制旁瓣};
    \draw[arrow] (gaincheck) -- node[left,pos=0.5] {是} (farfield);
    
    % Stage 6: Thermal check
    \node[step=red, below=1cm of farfield] (thermal) {Step 6: 热分析\\估算$R_{th}$, 确保$\Delta T < 50\,\mathrm{K}$};
    \draw[arrow] (farfield) -- (thermal);
    
    % Stage 7: Fabrication tolerance
    \node[step=brown, below=1cm of thermal] (tolerance) {Step 7: 工艺容差分析\\Monte Carlo 模拟$\pm 5\,\mathrm{nm}$偏差};
    \draw[arrow] (thermal) -- (tolerance);
    
    % Final output
    \node[step=blue, minimum height=1.2cm, below=1cm of tolerance] (output) {Step 8: 输出 GDSII 版图\\准备流片文档};
    \draw[arrow] (tolerance) -- (output);
    
    % Side notes: Tools recommendation
    \node[draw, dashed, align=left, font=\footnotesize, anchor=north east] at (6,9) {
      \textbf{推荐工具链}:\\
      $\bullet$ 能带:MPB/MITPhotonic Bands\\
      $\bullet$ 模式:Lumerical MODE/FIMMWAVE\\
      $\bullet$ 时域:MEEP/Lumerical FDTD\\
      $\bullet$ 多物理:COMSOL/ANSYS
    };
    
    % Timeline estimate
    \node[align=center, font=\small, anchor=south west] at (-5,0) {
      \textbf{典型开发周期}:\\
      设计迭代:2--4 周\\
      流片 + 测试:8--12 周\\
      总计:3--4 个月
    };
  \end{tikzpicture}
  \caption{PCSEL 从指标设定、初始结构设计、PWEM 选模、阈值分析、远场与热优化，到工艺容差评估和版图输出的完整工程流程。图中同时标出关键决策节点、典型回溯路径以及各阶段常用的仿真工具链。}
  \label{fig:ch22_design_workflow}
\end{figure}
